\documentclass[../Main.tex]{subfiles}
\begin{document}
  \chapter{2019-2020学年微积分（一）（上）期末考试}

  \section{单项选择题(每小题 3 分，共 18 分)}
  \begin{enumerate}
    \item 设$0<a_{n}<1$，且$\lim_{n\to\infty}a_{n}=1$，则以下数列中无界的是（\qquad）.

      \begin{tasks}[label=\textbf{\Alph*.}, label-width=1.5em, column-sep=1em](2)
        \task $\left\{a_{n}^{2}\right\}$
        \task $\left\{\frac{1}{a_{n}}\right\}$
        \task $\left\{\tan\frac{\pi a_{n}}{2}\right\}$
        \task $\left\{\ln a_{n}\right\}$
      \end{tasks}
      \vspace{1em}

    \item 已知$f(2)=3$，$f'(2)=5$，则极限$\lim_{h\to0}\frac{f^{2}(2+h)-9}{h}=$（\qquad）.

      \begin{tasks}[label=\textbf{\Alph*.}, label-width=1.5em, column-sep=1em](2)
        \task $30$
        \task $10$
        \task $6$
        \task $0$
      \end{tasks}
      \vspace{1em}

    \item 设函数$f(x)$在闭区间$[a,b]$上可导，则以下说法中\textbf{错误}的是（\qquad）.

      \begin{tasks}[label=\textbf{\Alph*.}, label-width=1.5em, column-sep=1em](2)
        \task $f(x)$必在$[a,b]$上有界
        \task $f(x)$必在$[a,b]$上有连续的导数
        \task $f(x)$必在$[a,b]$上连续
        \task $f(x)$必在$[a,b]$上可积
      \end{tasks}
      \vspace{1em}

    \item 曲线$y=\frac{1}{x}+\ln\left(1+\mathrm{e}^{x}\right)$一共有（\qquad）条渐近线.

      \begin{tasks}[label=\textbf{\Alph*.}, label-width=1.5em, column-sep=1em](2)
        \task $0$
        \task $1$
        \task $2$
        \task $3$
      \end{tasks}
      \vspace{1em}

    \item 设$f(x)$二阶可导，且$f''(x)<0$，则以下不等式中一定成立的是（\qquad）.

      \begin{tasks}[label=\textbf{\Alph*.}, label-width=1.5em, column-sep=1em](2)
        \task $f(1)+f(3)>2f(2)$
        \task $f(1)+f(3)<2f(2)$
        \task $f(1)+f(2)>2f(3)$
        \task $f(1)+f(2)<2f(3)$
      \end{tasks}
      \vspace{1em}

    \item 微分方程$y'-\frac{y}{2x}=0$满足初值条件$y(1)=2$的特解为（\qquad）.

      \begin{tasks}[label=\textbf{\Alph*.}, label-width=1.5em, column-sep=1em](2)
        \task $2\sqrt{x}$
        \task $1+\sqrt{x}$
        \task $1+x$
        \task $\sqrt{x+3}$
      \end{tasks}
      \vspace{1em}
  \end{enumerate}

  \section{填空题(每小题 4 分，共 16 分)}
  \begin{enumerate}
    \setcounter{enumi}{6}
    \item 极限$\lim_{n\to\infty}\frac{1}{n^{4}}\left(1^{3}+2^{3}+\cdots+n^{3}
      \right)=\underline{\hspace{3cm}}$.
      \vspace{0.5em}

    \item 设函数$y=x^{2}-x$. 在$x=2$，$\Delta x=0.01$时，微分$\mathrm{d}y=\underline
      {\hspace{3cm}}$.
      \vspace{0.5em}

    \item $\lim_{x\to0}\frac{\displaystyle\int_{0}^{\sin^2 x}\sqrt{3+t^2}\,\mathrm{d}t}{x^{2}}
      =\underline{\hspace{3cm}}$.
      \vspace{0.5em}

    \item 曲线$y=1-x^{4}$与$x$轴所围成图形的面积为$\underline{\hspace{3cm}}$.
  \end{enumerate}

  \section{计算题(每小题 7 分，共 42 分)}
  \begin{enumerate}
    \setcounter{enumi}{10}
    \item 已知当$x\to0$时，$u=\mathrm{e}^{3x}-ax^{2}-(1+bx)\cos x$是与$x^{3}$同阶的无穷小.
      求常数$a,b$的值.
      \vspace{12em}

    \item 求$f(x)=x^{3}-3x^{2}-9x+5$在区间$[-2,2]$上的最大值与最小值.
      \vspace{12em}

    \item 求不定积分$I=\int \frac{\mathrm{e}^{5x}}{\mathrm{e}^{2x}+1}\,\mathrm{d}
      x$.
      \vspace{11em}

    \item 求定积分$I=\int_{0}^{\frac{\pi}{4}}\frac{x}{\cos^{2} x}\,\mathrm{d}x$.
      \vspace{11em}

    \item 判定反常积分$\int_{0}^{1}\frac{1}{x\sqrt{x^{2}+1}}\,\mathrm{d}x$的敛散性，若收敛求其值.
      \vspace{12em}

    \item 求微分方程$y''=\frac{\sin y}{\cos^{3} y}$满足初值条件$y(2)=0$，$y'
      (2)=1$的特解.
      \vspace{12em}
  \end{enumerate}

  \section{综合题(每小题 7 分，共 14 分)}
  \begin{enumerate}
    \setcounter{enumi}{16}
    \item 设函数$f(x)=
      \begin{cases}
        \frac{\ln(1+x)}{x} , & x\in(-1,0)\cup(0,+\infty), \\
        1,                   & x=0
      \end{cases}$，讨论$f'(x)$在点$x=0$处的连续性.
      \vspace{12em}

    \item 求曲线$y=\sqrt{x}-\frac{x\sqrt{x}}{3}$对应于$1\leq x\leq 4$弧段的长度.
      \vspace{15em}
  \end{enumerate}

  \section{证明题(每小题 5 分，共 10 分)}
  \begin{enumerate}
    \setcounter{enumi}{18}
    \item 设$n$为正整数，求方程$(x+1)^{2n}=x^{2n}+1$所有的实根.证明你的结论.
      \vspace{15em}

    \item 设函数$f(x)$在区间$[a,b]$上有连续的导数. 证明：
      \[
        f(x)\leq \frac{1}{b-a}\left|\int_{a}^{b}f(x)\,\mathrm{d}x\right|+\int_{a}^{b}
        \left|f'(x)\right|\mathrm{d}x,\quad x\in[a,b].
      \]
      \vspace{15em}
  \end{enumerate}

  \chapter{2019-2020学年微积分（一）（上）期末考试参考答案}
  \section{单项选择题 (每小题 3 分，共 18 分)}
  \begin{enumerate}
    \item \textbf{Solution}. C.

      当$n\to\infty$时，$a_{n}\to1$且$0<a_{n}<1$，所以$\tan\frac{\pi a_{n}}{2}\to
      +\infty$，数列$\left\{\tan\frac{\pi a_{n}}{2}\right\}$无界.

    \item \textbf{Solution}. A.
      \[
        \begin{aligned}
          \lim_{h\to0}\frac{f^2(2+h)-9}{h}= & \lim_{h\to0}\frac{f^2(2+h)-f^2(2)}{h}=\lim_{h\to0}\frac{\left[f(2+h)-f(2)\right]\left[f(2+h)+f(2)\right]}{h} \\
          =                                 & \lim_{h\to0}\frac{f(2+h)-f(2)}{h}\cdot \lim_{h\to0}\left[f(2+h)+f(2)\right]                                  \\
          =                                 & 2f(2)f'(2)=30.
        \end{aligned}
      \]

    \item \textbf{Solution}. B.

      $f(x)$在闭区间$[a,b]$上必然有界，可导则说明$f(x)$在$[a,b]$上连续，必然可积.

      但可导不一定有连续的导数，故B选项错误.

    \item \textbf{Solution}. D.

      当$x\to0$时，$\lim_{x\to0}\left[\frac{1}{x}+\ln\left(1+\mathrm{e}^{x}\right
      )\right]=\infty$，所以曲线$y=\frac{1}{x}+\ln\left(1+\mathrm{e}^{x}\right)$有竖直渐近线$x
      =0$.

      当$x\to+\infty$时，$\lim_{x\to+\infty}\frac{\frac{1}{x}+\ln\left(1+\mathrm{e}^{x}\right)}{x}
      =\lim_{x\to+\infty}\frac{\ln\left(1+\mathrm{e}^{x}\right)}{x}=1$，

      $\lim_{x\to+\infty}\left[\frac{1}{x}+\ln\left(1+\mathrm{e}^{x}\right)-x\right
      ]=\lim_{x\to+\infty}\ln\frac{1+\mathrm{e}^{x}}{\mathrm{e}^{x}}=\lim_{x\to+\infty}
      \ln\left(1+\mathrm{e}^{-x}\right)=0$，

      所以曲线$y=\frac{1}{x}+\ln\left(1+\mathrm{e}^{x}\right)$有斜渐近线$y=x$.

      当$x\to-\infty$时，$\lim_{x\to-\infty}\frac{\frac{1}{x}+\ln\left(1+\mathrm{e}^{x}\right)}{x}
      =0$，$\lim_{x\to-\infty}\left[\frac{1}{x}+\ln\left(1+\mathrm{e}^{x}\right)\right
      ]=0$，

      所以曲线$y=\frac{1}{x}+\ln\left(1+\mathrm{e}^{x}\right)$有水平渐近线$y=0$.

      故曲线$y=\frac{1}{x}+\ln\left(1+\mathrm{e}^{x}\right)$的渐近线条数为3.

    \item \textbf{Solution}. B.

      $f(x)$是上凸的，所以$\frac{f(1)+f(3)}{2}<f(2)$，即$f(1)+f(3)<2f(2)$.

    \item \textbf{Solution}. A.

      方程变形为$\frac{\mathrm{d}y}{y}=\frac{\mathrm{d}x}{2x}$，两边积分得$\ln y=
      \frac{1}{2}\ln x+\ln C$，即$y=C\sqrt{x}$.

      将初值条件$y(1)=2$代入上式得$C=2$，所以特解为$y=2\sqrt{x}$.
  \end{enumerate}

  \section{填空题(每小题 4 分，共 16 分)}
  \begin{enumerate}
    \setcounter{enumi}{6}
    \item \textbf{Solution}. $\frac{1}{4}$.

      由定积分的定义，
      \[
        \begin{aligned}
          \lim_{n\to\infty}\frac{1}{n^4}\left(1^{3}+2^{3}+\cdots+n^{3}\right) = & \lim_{n\to\infty}\frac{1}{n}\left(\frac{1^3}{n^3}+\frac{2^3}{n^3}+\cdots+\frac{n^3}{n^3}\right) \\
          =                                                                     & \lim_{n\to\infty}\frac{1}{n}\sum_{i=1}^{n}\left(\frac{i}{n}\right)^{3}                          \\
          =                                                                     & \int_{0}^{1}x^{3}\,\mathrm{d}x = \frac{1}{4}.
        \end{aligned}
      \]

    \item \textbf{Solution}. $0.03$.

      $\mathrm{d}y=y'(2)\Delta x = (2x-1)\bigg|_{x=2}\Delta x=3\Delta x=0.03$.

    \item \textbf{Solution}. $\sqrt{3}$.
      \[
        \begin{aligned}
          \lim_{x\to0}\frac{\displaystyle\int_{0}^{\sin^2 x}\sqrt{3+t^2}\,\mathrm{d}t}{x^2}= & \lim_{x\to0}\frac{\sqrt{3+\sin ^4 x}\cdot 2\sin x\cos x}{2x} \\
          =                                                                                & \sqrt{3}\lim_{x\to0}\cos x = \sqrt{3}.
        \end{aligned}
      \]

    \item \textbf{Solution}. $\frac{8}{5}$.

      所围成图形的面积$S=\int_{-1}^{1}\left(1-x^{4}\right)\,\mathrm{d}x = \frac{8}{5}$.
  \end{enumerate}
  \section{计算题(每小题 7 分，共 42 分)}
  \begin{enumerate}
    \setcounter{enumi}{10}
    \item \textbf{Solution}. 由Taylor公式，
      \[
        \begin{aligned}
          u = & \mathrm{e}^{3x}-ax^{2}-(1+bx)\cos x                                                                                \\
          =   & \left(1+3x+\frac{9}{2}x^{2}+\frac{9}{2}x^{3}+o(x^{3})\right)-ax^{2}-(1+bx)\left(1-\frac{1}{2}x^{2}+o(x^{3})\right) \\
          =   & (3-b)x+(5-a)x^{2}+\frac{9+b}{2}x^{3}+o(x^{3}).
        \end{aligned}
      \]

      因此必然有$a=5$，$b=3$.

    \item \textbf{Solution}. $f'(x)=3x^{2}-6x-9$，令$f'(x)=0$得到函数在$(-2,
      2)$内的驻点为$x=-1$.

      计算$f(-2)=3$，$f(-1)=10$，$f(2)=-17$，所以函数在区间$[-2,2]$上的最大值为$1
      0$，最小值为$-17$.

    \item \textbf{Solution}. 令$u=\mathrm{e}^{x}$，则$\mathrm{d}u=\mathrm{e}
      ^{x}\,\mathrm{d}x$，
      \[
        \begin{aligned}
          I = & \int \frac{\mathrm{e}^{4x}\cdot\mathrm{e}^x}{\left(\mathrm{e}^x\right)^2+1}\,\mathrm{d}x = \int \frac{u^4}{u^2+1}\,\mathrm{d}u \\
          =   & \int \left(u^{2}-1+\frac{1}{u^2+1}\right)\,\mathrm{d}u                                                                       \\
          =   & \frac{u^3}{3}-u+\arctan u + C = \frac{\mathrm{e}^{3x}}{3}-\mathrm{e}^{x}+\arctan\mathrm{e}^{x}+C.
        \end{aligned}
      \]

    \item \textbf{Solution}.
      \[
        \begin{aligned}
          I = \int_{0}^{\frac{\pi}{4}}x\,\mathrm{d}\tan x = & x\tan x\bigg|_{0}^{\frac{\pi}{4}}-\int_{0}^{\frac{\pi}{4}}\tan x\,\mathrm{d}x \\
          =                                               & \frac{\pi}{4}+\ln|\cos x|\bigg|_{0}^{\frac{\pi}{4}}                         \\
          =                                               & \frac{\pi}{4}-\frac{1}{2}\ln 2.
        \end{aligned}
      \]

    \item \textbf{Solution}. 令$x=\tan t$，则$\mathrm{d}x=\sec^{2} t\,\mathrm{d}
      t$，
      \[
        \begin{aligned}
          \int_{0}^{1}\frac{1}{x\sqrt{x^2+1}}\,\mathrm{d}x = & \int_{0}^{\frac{\pi}{4}}\frac{\sec^2 t}{\tan t\sqrt{\tan^2 t+1}}\,\mathrm{d}t = \int_{0}^{\frac{\pi}{4}}\csc t\,\mathrm{d}t \\
          =                                                & -\ln|\csc t+\cot t|\bigg|_{0^+}^{\frac{\pi}{4}}= -\ln(\sqrt{2}+1)+\lim_{t\to0^+}\ln\frac{1+\cos t}{\sin t}              \\
          =                                                & +\infty.
        \end{aligned}
      \]

      所以反常积分$\int_{0}^{1}\frac{1}{x\sqrt{x^{2}+1}}\,\mathrm{d}x$发散.

    \item \textbf{Solution}. 令$y'=p$，则$y''=\frac{\mathrm{d}p}{\mathrm{d}x}
      =\frac{\mathrm{d}p}{\mathrm{d}y}\cdot \frac{\mathrm{d}y}{\mathrm{d}x}=p\frac{\mathrm{d}p}{\mathrm{d}y}$，

      微分方程$y''=\frac{\sin y}{\cos^{3} y}$变形为$p\,\mathrm{d}p=\frac{\sin y}{\cos^{3}
      y}\,\mathrm{d}y$.

      方程两边积分得$\frac{1}{2}p^{2}=\frac{1}{2}\sec^{2} y+C_{1}$，将初值条件$y(
      2)=0$，$y'(2)=1$代入上式得$C_{1}=0$，

      所以$p=\sec y$，即$\frac{\mathrm{d}y}{\mathrm{d}x}=\sec y$或$\cos y\,\mathrm{d}
      y=\mathrm{d}x$.

      方程两边积分得$x=\sin y+C_{2}$，将初值条件$y(2)=0$代入上式得$C_{2}=2$，

      所以特解为$x=\sin y+2$或$y=\arcsin(x-2)$.
  \end{enumerate}

  \section{综合题(每小题 7 分，共 14 分)}
  \begin{enumerate}
    \setcounter{enumi}{16}
    \item \textbf{Solution}. 当$x\in(-1,0)\cup(0,+\infty)$时，$f'(x)=\frac{\frac{x}{1+x}-\ln(1+x)}{x^{2}}
      =\frac{x-(1+x)\ln(1+x)}{x^{2}(1+x)}$.

      当$x=0$时，$f'(0)=\lim_{x\to0}\frac{f(x)-f(0)}{x}=\lim_{x\to0}\frac{\ln(1+x)-x}{x^{2}}
      =-\frac{1}{2}$.

      又$\lim_{x\to0}\frac{x-(1+x)\ln(1+x)}{x^{2}(1+x)}=\lim_{x\to0}\frac{x-(1+x)\ln(1+x)}{x^{2}}
      =\lim_{x\to0}\frac{1-\ln(1+x)-1}{2x}=-\frac{1}{2}=f'(0)$，

      所以$f'(x)$在点$x=0$处连续.

    \item \textbf{Solution}.$y'=\frac{1}{2\sqrt{x}}-\frac{\sqrt{x}}{2}$.

      所以$\mathrm{d}s=\sqrt{1+y'^{2}}\,\mathrm{d}x=\sqrt{1+\frac{1}{4x}+\frac{x}{4}-\frac{1}{2}}
      \mathrm{d}x=\sqrt{\left(\frac{1}{2\sqrt{x}}+\frac{\sqrt{x}}{2}\right)^{2}}\,\mathrm{d}
      x=\frac{1}{2}\left(\frac{1}{\sqrt{x}}+\sqrt{x}\right)\,\mathrm{d}x$.

      故$s=\frac{1}{2}\int_{1}^{4}\left(\frac{1}{\sqrt{x}}+\sqrt{x}\right)\,\mathrm{d}
      x=\frac{1}{2}\left[2\sqrt{x}+\frac{2}{3}x^{\frac{3}{2}}\right]_{1}^{4}=\frac{1}{2}
      \left(4+\frac{16}{3}-2-\frac{2}{3}\right)=\frac{10}{3}$.
  \end{enumerate}

  \section{证明题(每小题 5 分，共 10 分)}
  \begin{enumerate}
    \setcounter{enumi}{18}
    \item \textbf{Proof}. 显然$x=0$是方程$(x+1)^{2n}=x^{2n}+1$的一个实根.

      令$f(x)=(x+1)^{2n}-x^{2n}-1$，则$f'(x)=2n\left[(x+1)^{2n-1}-x^{2n-1}\right]$.

      由于$t^{2n-1}$是严格单增函数，所以$f'(x)>0$，$f(x)$在$(-\infty,+\infty)$上严格单增.

      因此方程$(x+1)^{2n}=x^{2n}+1$只有唯一实根$x=0$.

    \item \textbf{Solution}. 根据积分中值定理，$\exists\,\xi\in[a,b]$，使得$f
      (\xi)=\frac{1}{b-a}\int_{a}^{b}f(x)\,\mathrm{d}x$.

      所以
      \[
        \begin{aligned}
          f(x) = & f(\xi)+f(x)-f(\xi) \leq |f(\xi)|+\left|\int_{\xi}^{x}f'(t)\,\mathrm{d}t\right|                          \\
          \leq   & \frac{1}{b-a}\left|\int_{a}^{b}f(x)\,\mathrm{d}x\right| + \left|\int_{\xi}^{x}|f'(t)|\mathrm{d}t\right| \\
          \leq   & \frac{1}{b-a}\left|\int_{a}^{b}f(x)\,\mathrm{d}x\right| + \int_{a}^{b}|f'(t)|\mathrm{d}t.
        \end{aligned}
      \]
  \end{enumerate}
\end{document}




